% I. PLANTEAMIENTO DEL PROBLEMA. Lo escribe la skill tesis-problema.
% El título del capítulo lo pone generado/estructura.tex; aquí van las secciones.

\section{Descripción del problema}
% Realidad problemática: internacional -> nacional -> local/organización.
% Síntomas -> causas -> pronóstico -> control al pronóstico. Toda cifra o afirmación con cita.
% Sin subtítulos, en tercera persona; incluye motivación y estado del arte del problema.

\section{Formulación del problema}
% GENERADO por generar.py desde tesis.yaml. No editar a mano.
\subsection{Problema general}
¿En qué medida un sistema web con chatbot de WhatsApp e inteligencia artificial mejora la gestión de atención al cliente en una empresa de servicios?

\subsection{Problemas específicos}
\begin{enumerate}[label=\textbf{PE\arabic*:},leftmargin=1.27cm,itemsep=4pt]
\item ¿En qué medida un sistema web con chatbot de WhatsApp e inteligencia artificial reduce el tiempo promedio de primera respuesta al cliente?
\item ¿En qué medida un sistema web con chatbot de WhatsApp e inteligencia artificial incrementa la tasa de consultas resueltas sin intervención humana?
\end{enumerate}


\section{Objetivos de la investigación}
% GENERADO por generar.py desde tesis.yaml. No editar a mano.
\subsection{Objetivo general}
Determinar en qué medida un sistema web con chatbot de WhatsApp e inteligencia artificial mejora la gestión de atención al cliente en una empresa de servicios.

\subsection{Objetivos específicos}
\begin{enumerate}[label=\textbf{OE\arabic*:},leftmargin=1.27cm,itemsep=4pt]
\item Determinar en qué medida un sistema web con chatbot de WhatsApp e inteligencia artificial reduce el tiempo promedio de primera respuesta al cliente.
\item Determinar en qué medida un sistema web con chatbot de WhatsApp e inteligencia artificial incrementa la tasa de consultas resueltas sin intervención humana.
\end{enumerate}


\section{Delimitación de la investigación}
\subsection{Delimitación espacial}
\DelimEspacial

\subsection{Delimitación temporal}
\DelimTemporal

\section{Justificación del problema}
% Un \subsection por tipo, según corresponda: teórica, tecnológica, social, económica.
% Con cifras citadas.
